
\medskip
\phantomsection\label{cor:sga3-VIB-7.1.1}%
\noindent\textbf{Corollary 7.1.1.---}%
\footnote{\label{note:sga3-VIB-70}Editor's note: this corollary
has been added to point out this special case of \SGARefLink{prop:sga3-VIB-7.1}{7.1}.}
\emph{Let \(G\) be a \(k\)-group locally of finite type, let
\(A\) and \(B\) be two \(k\)-subgroups of \(G\) that are smooth
and of finite type, and let \(i_A\) (respectively, \(i_B\)) be the
inclusion of \(A\) (respectively, \(B\)) into \(G\). Suppose that
\(B\) normalizes \(A\). Then}
\phantomsection\label{form:sga3-VIB-7.1.1-generated-product}%
\[
  A\cdot B=\Gamma_G(i_A,i_B).
\]

\emph{Proof.}
Let
\[
  H=A\rtimes B
\]
be the semidirect product of \(A\) and \(B\) (cf.\ I, \SGARefLink{sga3:i:definition:2-3-5}{2.3.5});
it is a smooth \(k\)-group of finite type. The group morphism
\[
  u:H\longrightarrow G,
  \qquad
  (a,b)\longmapsto ab,
\]
is therefore quasi-compact. Hence, by
\hyperref[prop:sga3-VIB-1.2]{1.2},
\[
  u(H)=A\cdot B
\]
is a closed reduced subscheme of \(G\), and is a group in the
category \((\mathrm{Sch}/k)_{\mathrm{red}}\). By
EGA~IV\(_3\), 11.10.7, \(A\cdot B=u(H)\) is geometrically
reduced, since \(H\) is; it is therefore a closed subgroup of
\(G\). The evident inclusion
\[
  A\cdot B\subset\Gamma_G(i_A,i_B)
\]
now proves the corollary.

\par
\ifdim\dimexpr\pagegoal-\pagetotal\relax<12\baselineskip
  \newpage
\fi
\medskip
\phantomsection\label{defrem:sga3-VIB-7.2}%
\noindent\textbf{Definitions and \SGARefLink{defrem:sga3-VIB-7.2}{Remarks 7.2}.---}%
\footnote{\label{note:sga3-VIB-71}Editor's note: item
\textup{(viii)} of the original has been placed in
\textup{(i\(^{\prime}\))}, and items \textup{(vi)} and
\textup{(vii)} have been singled out as \hyperref[cor:sga3-VIB-7.2.1]{Corollary 7.2.1} and
\hyperref[def:sga3-VIB-7.2.2]{Definition 7.2.2} below.}
\begin{enumerate}
\item[\textup{(i)}]\phantomsection\label{def:sga3-VIB-7.2-i}
  Let \(G\) be a \(k\)-group, let \((X_i)_{i\in I}\) be a
  family of geometrically reduced \(k\)-schemes, and, for every
  \(i\in I\), let \(f_i:X_i\to G\) be a \(k\)-morphism. The
  \emph{closed group subscheme of \(G\) generated by the family
  \((f_i)_{i\in I}\)} is the smallest closed group subscheme of
  \(G\) through which every \(f_i\) factors; in this subsection
  it is denoted
  \[
    \Gamma_G\bigl((f_i)_{i\in I}\bigr).
  \]
  If \(X\) is a subscheme of \(G\), geometrically reduced over
  \(k\), and \(f:X\hookrightarrow G\) is the immersion, we write
  \(\Gamma_G(X)\) instead of \(\Gamma_G(f)\).

\item[\textup{(i\(^{\prime}\))}]
  \phantomsection\label{def:sga3-VIB-7.2-iprime}
  With the notation of
  \hyperref[concl:sga3-VIB-7.1-ii]{7.1\textup{(ii)}}, we shall
  sometimes put
  \phantomsection\label{form:sga3-VIB-7.2-iprime-Gamma}%
  \[
    \Gamma'_G(f)=\bigcup_{n\geq1}f^n(X^n).
  \]
  Notice that \(\Gamma'_G(f)\) is a subset of \(G\) stable
  under the group law, in the sense of \hyperref[par:sga3-VIB-3.0]{3.0}.

\item[\textup{(ii)}]\phantomsection\label{rem:sga3-VIB-7.2-ii}
  If \(X_1\) and \(X_2\) are geometrically reduced
  \(k\)-schemes and \(f_1:X_1\to G\) and \(f_2:X_2\to G\) are
  \(k\)-morphisms such that
  \[
    f_1(X_1)=f_2(X_2)
  \]
  as sets, then it is clear that
  \[
    \Gamma_G(f_1)=\Gamma_G(f_2).
  \]

\item[\textup{(iii)}]\phantomsection\label{def:sga3-VIB-7.2-iii}
  Let \(E\) be a subset of \(G\) such that the reduced
  subscheme \(\overline E\) of \(G\) is geometrically reduced.
  The \emph{closed group subscheme of \(G\) generated by \(E\)}
  is the group subscheme
  \[
    \Gamma_G(E):=\Gamma_G(i),
  \]
  where \(i:\overline E\hookrightarrow G\) is the inclusion of
  that reduced subscheme.

\item[\textup{(iv)}]\phantomsection\label{rem:sga3-VIB-7.2-iv}
  Since every group subscheme of \(G\) is closed, by
  \hyperref[sga3:VIA:0.5.2]{VI\(_A\), 0.5.2},%
  we shall speak of the ``generated group subscheme'' rather
  than the ``generated closed group subscheme.''

\item[\textup{(v)}]\phantomsection\label{rem:sga3-VIB-7.2-v}
  Let \(X\) be a geometrically reduced \(k\)-scheme and let
  \(f:X\to G\) be a \(k\)-morphism. Suppose that \(f(X)\)
  contains the unit element \(e\) of \(G\). Put
  \[
    X'^{\,1}=X\times_kX,
    \qquad
    f'^{\,1}=\mu\circ\bigl(f\times_k(c\circ f)\bigr),
  \]
  and, for \(n>1\), put
  \phantomsection\label{form:sga3-VIB-7.2-v-recursion}%
  \[
    X'^{\,n}=X'^{\,1}\times_kX'^{\,n-1},
    \qquad
    f'^{\,n}
      =\mu\circ\bigl(f'^{\,1}\times_k f'^{\,n-1}\bigr).
  \]
  Then \(\Gamma_G(f)\) is the reduced subscheme of \(G\) whose
  underlying space is the closure of
  \[
    \bigcup_{n\geq1}f'^{\,n}(X'^{\,n}).
  \]

  Indeed, recall the notation of
  \hyperref[concl:sga3-VIB-7.1-ii]{7.1\textup{(ii)}}:
  \[
    X^1=X\amalg X,
    \qquad
    X^n=X^1\times_kX^{n-1}\quad(n\geq2).
  \]
  The hypothesis \(e\in f(X)\) gives, for every \(n\geq1\),
  \phantomsection\label{form:sga3-VIB-7.2-v-inclusions}%
  \[
    f^n(X^n)
      \subset f'^{\,n}(X'^{\,n})
      \subset f^{2n}(X^{2n}).
  \]
  Consequently, there exists an integer \(n\) such that
  \(f^n(X^n)=\Gamma_G(f)\) if and only if there exists an
  integer \(m\) such that
  \(f'^{\,m}(X'^{\,m})=\Gamma_G(f)\).
\end{enumerate}

From \hyperref[rem:sga3-VIB-7.2-v]{Remark 7.2\textup{(v)}} we
deduce the following result.

\medskip
\phantomsection\label{cor:sga3-VIB-7.2.1}%
\noindent\textbf{Corollary 7.2.1.---}
\emph{Let \(X\) be a geometrically reduced and geometrically
connected \(k\)-scheme, and let \(f:X\to G\) be a \(k\)-morphism
such that \(f(X)\) contains the unit element of \(G\). Then the
\(k\)-group \(\Gamma_G(f)\) is connected, and hence
irreducible.}

\emph{Proof.}
Each \(X'^{\,n}\) is connected, so the union of the
\[
  f'^{\,n}(X'^{\,n}),
\]
all of which contain the unit element, is connected. Its closure
\(\Gamma_G(f)\) is therefore connected as well. It follows that
\(\Gamma_G(f)\) is irreducible, by
\hyperref[prop:sga3-VIA-2.4]{VI\(_A\), 2.4} when \(G\), and
hence also \(\Gamma_G(f)\), is locally of finite type over
\(k\), and by
\hyperref[thm:sga3-VIA-2.6.5]{VI\(_A\), 2.6.5\textup{(iii)}}
in the general case.
